\chapter{Results, Discussion and Perspectives}

\section*{Introduction}

This chapter compares the results obtained against the initial objectives, analyses
the platform's strengths and limitations objectively, and describes the evolution
perspectives identified during development.

\section{Results Obtained}

\begin{table}[H]
\centering
\caption{Objectives vs.\ results matrix}
\label{tab:obj-results}
\begin{tabularx}{\textwidth}{lXl}
\toprule
\textbf{Objective} & \textbf{Result} & \textbf{Status} \\
\midrule
O1: Answers-first dashboard
  & Status banner, issues, recommended actions
  & \textcolor{SuccessGreen}{\textbf{OK}} \\
O2: Change correlation
  & Full engine: temporal scoring + LLM narrative
  & \textcolor{SuccessGreen}{\textbf{OK}} \\
O3: Local LLM explanation
  & Ollama + Llama 3.1 8B, non-DevOps prompts
  & \textcolor{SuccessGreen}{\textbf{OK}} \\
O4: Holt-Winters 24h forecast
  & 24 points, time-to-critical, LLM narrative
  & \textcolor{SuccessGreen}{\textbf{OK}} \\
O5: Microservices architecture
  & 19 services, NATS JetStream, 8 streams
  & \textcolor{SuccessGreen}{\textbf{OK}} \\
O6: End-to-end security
  & mTLS, JWT, RLS, AES-256-GCM, audit log
  & \textcolor{SuccessGreen}{\textbf{OK}} \\
O7: Deployment $\leq$10 min
  & 6 steps, $\approx$8 min (excl.\ model download)
  & \textcolor{SuccessGreen}{\textbf{OK}} \\
\bottomrule
\end{tabularx}
\end{table}

\subsection{Codebase Metrics}

\begin{table}[H]
\centering
\caption{Quantitative codebase indicators}
\label{tab:codebase-metrics}
\begin{tabularx}{\textwidth}{lX}
\toprule
\textbf{Indicator} & \textbf{Value} \\
\midrule
Total lines of code (excl. tests and comments) & $\sim$52,000 \\
Python services   & 11 \\
Go services       & 1 (telemetry-writer) \\
React/TypeScript components & 26 components, 6 pages \\
SQL migrations    & 14 \\
PostgreSQL tables & 22 \\
NATS JetStream streams & 8 \\
Unit tests        & 34 pytest tests \\
Critical test coverage & $\sim$74\% (detectors + forecaster + scoring) \\
\bottomrule
\end{tabularx}
\end{table}

\section{Strengths}

\begin{description}[leftmargin=2cm]
  \item[Decoupled architecture] NATS JetStream ensures no service directly depends
    on another. A new anomaly consumer (e.g. an automated remediation service)
    can be added without modifying existing services.
  \item[Targeted natural-language explanation] LLM prompts written for a
    non-DevOps audience produce explanations directly understandable by
    developers managing their own servers.
  \item[Change correlation] The automatic change correlation engine is absent
    from all compared open source solutions. By linking anomalies to recent
    changes, it reduces diagnosis time from tens of minutes to seconds.
  \item[Zero cloud dependency] All AI processing runs locally. Data never leaves
    the operator's server, satisfying the strictest confidentiality requirements.
  \item[Portability] Docker Compose deployment works identically on Linux,
    macOS and Windows with Docker Desktop.
\end{description}

\section{Identified Limitations}

\begin{description}[leftmargin=2cm]
  \item[No seasonal component in Holt-Winters] The double exponential model
    does not capture cyclical variations (daily CPU peaks, weekly traffic).
    Triple Holt-Winters would improve 24h forecast accuracy.
  \item[LLM quality depends on hardware] On a machine without GPU, Llama 3.1 8B
    inference can take 15--30 seconds. Acceptable for asynchronous explanation
    but perceptible to the user.
  \item[Linux-only agent change reporters] The \texttt{auth.log}, \texttt{dpkg.log}
    and \texttt{syslog} watchers are Linux-specific. Windows server monitoring
    requires a separate reporter implementation.
  \item[No automated remediation] \caimp{} produces action recommendations but
    executes no corrections automatically. This is intentional (precautionary
    principle) but limits appeal for teams seeking autonomous operation.
\end{description}

\section{Future Perspectives}

\textbf{Short term (3--6 months):}
Triple Holt-Winters (seasonal component); Windows change reporters; native
Slack/Teams alerts; LLM fine-tuning on organisational incidents using
thumbs-up/down feedback dataset.

\textbf{Medium term (6--18 months):}
Trace correlation (OpenTelemetry traces in the correlation engine); assisted
automated remediation with operator confirmation; hosted SaaS offering
(Kubernetes per-tenant isolation); React Native mobile app with push notifications.

\textbf{Long-term vision:}
\caimp{} should become the virtual SRE of small teams — a system capable not
only of explaining past incidents but of autonomously preventing future ones,
while keeping the operator informed and in control.

\section*{Conclusion}

The results demonstrate that all objectives were fully achieved within the
one-month timeframe. \caimp{} delivers a unique value proposition in the open
source monitoring landscape: natural-language explanations, automatic change
correlation, and deployment accessible to all. The identified limitations form
a clear development roadmap for future versions.
